\documentclass[11pt,a4paper]{article}

% ========================================
% PACKAGES
% ========================================
\usepackage[utf8]{inputenc}
\usepackage[T1]{fontenc}
\usepackage{lmodern}
\usepackage[english]{babel}
\usepackage[margin=2.5cm]{geometry}

% Mathematics
\usepackage{amsmath, amssymb, amsthm}
\usepackage{mathtools}

% Graphics and colors
\usepackage{graphicx}
\usepackage{xcolor}
\usepackage{tikz}
\usetikzlibrary{positioning, arrows.meta, shapes.geometric, calc}
\usepackage{pgfplots}
\pgfplotsset{compat=1.18}
\usepackage[most]{tcolorbox}

% Lists and formatting
\usepackage{enumitem}
\usepackage{parskip}
\usepackage{fancyhdr}
\usepackage{multirow}

% Code listings
\usepackage{listings}
\usepackage{algorithm}
\usepackage{algorithmic}

% Hyperlinks
\usepackage{hyperref}
\hypersetup{
    colorlinks=true,
    linkcolor=blue,
    urlcolor=blue,
    citecolor=blue
}

% ========================================
% THEOREM ENVIRONMENTS
% ========================================
\theoremstyle{definition}
\newtheorem{definition}{Definition}[section]
\newtheorem{example}{Example}[section]
\newtheorem{exercise}{Exercise}[section]

\theoremstyle{plain}
\newtheorem{theorem}{Theorem}[section]
\newtheorem{lemma}[theorem]{Lemma}
\newtheorem{proposition}[theorem]{Proposition}
\newtheorem{corollary}[theorem]{Corollary}

\theoremstyle{remark}
\newtheorem*{remark}{Remark}
\newtheorem*{observation}{Observation}

% ========================================
% CUSTOM COMMANDS
% ========================================
\newcommand{\R}{\mathbb{R}}
\newcommand{\N}{\mathbb{N}}
\newcommand{\Z}{\mathbb{Z}}
\newcommand{\Q}{\mathbb{Q}}
\newcommand{\C}{\mathbb{C}}

% Robotics-specific
\newcommand{\SO}[1]{\mathrm{SO}(#1)} % chktex 36
\newcommand{\SE}[1]{\mathrm{SE}(#1)} % chktex 36
\newcommand{\T}[2]{{}^{#1}T_{#2}}          % homogeneous transform: ^A T_B
\newcommand{\p}[1]{\mathbf{#1}}             % bold vector
\newcommand{\skewmat}[1]{[#1]_{\times}}     % skew-symmetric matrix
\newcommand{\norm}[1]{\left\lVert#1\right\rVert}

% ========================================
% TABLE OF CONTENTS DEPTH
% ========================================
\setcounter{tocdepth}{2}

% ========================================
% HEADER AND FOOTER
% ========================================
\setlength{\headheight}{14pt}
\pagestyle{fancy}
\fancyhf{}
\lhead{\leftmark}
\rhead{AI Robotics}
\cfoot{\thepage}
\renewcommand{\sectionmark}[1]{\markboth{#1}{}}

% ========================================
% DOCUMENT INFORMATION
% ========================================
\title{\textbf{AI Robotics}\\
\large Artificial Intelligence}
\author{Jacopo Parretti}
\date{II Semester 2025--2026}

% ========================================
% DOCUMENT
% ========================================
\begin{document}

\maketitle
\tableofcontents
\newpage


% ========================================
\part{Mobile Robotics: Introduction}
% ========================================

\section{Autonomous Mobile Robots}

\subsection{The Three Key Questions}

Mobile robotics revolves around three fundamental questions that any autonomous system must answer:

\begin{enumerate}
    \item \textbf{Where am I?} $\to$ Perception + Localization + Mapping
    \item \textbf{Where am I going?} $\to$ Planning + Path Planning
    \item \textbf{How do I get there?} $\to$ Obstacle Avoidance + Kinematics
\end{enumerate}

To answer these questions, the robot must:
\begin{itemize}
    \item have a \textbf{model} of the environment
    \item \textbf{perceive} and analyze the environment
    \item find its \textbf{position/situation} in the environment
    \item \textbf{plan} and execute the movement
\end{itemize}

\subsection{General Control Scheme}

The general architecture of an autonomous mobile robot can be decomposed into two main pipelines:

\begin{itemize}
    \item \textbf{Perception pipeline}: Sensing $\to$ Information Extraction and Interpretation $\to$ Localization and Map Building
    \item \textbf{Motion control pipeline}: Cognition and Path Planning $\to$ Path Execution $\to$ Acting
\end{itemize}

The \textit{Localization/Map Building} block feeds a position estimate and a global map to the \textit{Cognition} block, which uses mission commands from the knowledge base to plan a path. The path is then passed to \textit{Path Execution}, which sends actuator commands to the \textit{Acting} block that interacts with the real-world environment. The loop is closed by sensing the resulting state.

\section{Kinematics and Motion Control}

\subsection{Wheel Types and Constraints}

Wheeled mobile robots are subject to two fundamental kinematic constraints:

\begin{itemize}
    \item \textbf{Rolling constraint}: the wheel rolls without slipping in the rolling direction — velocity component along the wheel plane must be consistent with wheel rotation.
    \item \textbf{No-sliding constraint}: the wheel does not slip laterally — zero velocity component orthogonal to the wheel plane.
\end{itemize}

Different wheel types (fixed standard, steerable, castor, Swedish/Mecanum) impose different combinations of these constraints and determine the robot's degrees of freedom and maneuverability.

\subsection{Forward and Inverse Kinematics}

Let $(x, y, \theta)$ be the robot pose in the global frame and $\varphi_1, \ldots, \varphi_n$ the wheel rotation speeds. The \textbf{forward kinematics} (wheels $\to$ body velocity) and \textbf{inverse kinematics} (body velocity $\to$ wheels) are:

\[
\begin{bmatrix} \dot{x} \\ \dot{y} \\ \dot{\theta} \end{bmatrix}
= f(\dot{\varphi}_1, \ldots, \dot{\varphi}_n, \theta, \text{geometry})
\qquad
\begin{bmatrix} \dot{\varphi}_1 \\ \vdots \\ \dot{\varphi}_n \end{bmatrix}
= f(\dot{x}, \dot{y}, \dot{\theta})
\]

The geometry of the robot (wheel placement, radius, orientation) determines the specific form of $f$.

\section{Perception}

\subsection{Sensing Modalities}

Perception is the process of extracting meaningful information from raw sensor data. Common sensors on mobile robots include:

\begin{itemize}
    \item \textbf{Lidar} (2D or 3D): provides range measurements by emitting laser pulses and measuring time-of-flight. Produces dense point clouds or scan profiles used for obstacle detection, mapping, and localization.
    \item \textbf{Cameras}:
    \begin{itemize}
        \item \textit{Stereo}: two horizontally offset cameras enabling triangulation-based depth estimation.
        \item \textit{Omnidirectional}: wide field-of-view imaging via catadioptric optics.
        \item \textit{RGB-D}: combines color image with per-pixel depth (structured light or ToF).
    \end{itemize}
\end{itemize}

Laser scans are frequently overlaid on occupancy grid maps for real-time spatial reasoning.

\section{Localization and Mapping}

\subsection{Localization: The Core Problem}

Localization is the task of estimating the robot's pose $(x, y, \theta)$ relative to a known or partially known map. It is inherently probabilistic: sensors are noisy, and motion estimation via encoders accumulates error over time.

\subsection{Markov Localization}

A classical probabilistic approach to localization proceeds through a \textbf{sense–act–sense} cycle:

\begin{enumerate}
    \item \textbf{SEE}: robot queries its sensors $\to$ observes a landmark (e.g.\ a door). The belief distribution is updated (sharpened at consistent hypotheses).
    \item \textbf{ACT}: robot moves forward. Motion is estimated via encoders; uncertainty grows (belief spreads).
    \item \textbf{SEE}: robot observes again. The new observation is fused with the propagated belief via Bayes' rule.
\end{enumerate}

\begin{tcolorbox}[colback=blue!5!white,colframe=blue!75!black,title=\textbf{Belief Update}]
At each step, the robot maintains a probability distribution over poses — the \textbf{belief} $\text{bel}(x_t)$. It is updated by:
\[
\text{bel}(x_t) = \eta \, p(z_t \mid x_t) \int p(x_t \mid u_t, x_{t-1})\, \text{bel}(x_{t-1})\, dx_{t-1}
\]
where $z_t$ is the sensor observation, $u_t$ is the control input, and $\eta$ is a normalizing constant.
\end{tcolorbox}

\subsection{Visual Odometry}

An alternative to encoder-based motion estimation is \textbf{visual odometry}: estimating ego-motion from consecutive camera frames. Modern methods such as \textit{Direct Sparse Odometry} (DSO) operate directly on pixel intensities without requiring explicit feature matching, achieving robust and accurate trajectory estimation.

\section{Cognition}

\subsection{Obstacle Avoidance}

Obstacle avoidance employs \textbf{local or reactive methods} to steer the robot toward its goal while avoiding collisions in real time. Two representative approaches:

\begin{itemize}
    \item \textbf{Bug algorithms} (e.g.\ Bug2): the robot moves in a straight line toward the goal and follows obstacle boundaries when a collision is imminent. Simple and provably complete for certain environments.
    \item \textbf{Vector Field Histogram (VFH)}: builds a polar histogram of obstacle density around the robot; selects a steering direction in a sector with low obstacle density aligned with the goal direction.
\end{itemize}

\subsection{Global Path Planning}

Global path planning computes a \textbf{collision-free path} from the robot's current position to the goal using a prior map. Representative methods:

\begin{itemize}
    \item \textbf{Voronoi diagram}: paths are planned along the medial axis of the free space, maximizing clearance from obstacles.
    \item \textbf{Cell decomposition}: the free space is partitioned into cells; a graph is built over adjacent free cells and a graph-search algorithm (e.g.\ A$^*$) finds the optimal path.
\end{itemize}

\subsection{Decision Making}

Beyond reactive and planned navigation, a robot may need to \textbf{regulate its behavior online} — e.g.\ adjust velocity to complete a path as fast as possible while avoiding accidents under uncertainty.

\textbf{Online Monte Carlo Planning} (in particular POMCP — Partially Observable Monte Carlo Planning) addresses this by:
\begin{itemize}
    \item modeling the task as a POMDP (Partially Observable Markov Decision Process)
    \item using Monte Carlo Tree Search to plan online without a full model of the environment
    \item incorporating \textbf{prior knowledge} and handling \textbf{state uncertainty} explicitly through belief states
\end{itemize}

\section{Learning}

\subsection{Deep Reinforcement Learning for Robotic Navigation}

When explicit maps are unavailable or the environment is non-stationary, \textbf{Deep Reinforcement Learning (DRL)} provides a powerful paradigm for learning navigation policies directly from experience.

\textbf{Mapless navigation for wheeled robots}: a DRL agent (e.g.\ DDQN with discrete action space) is trained in simulation (Unity) and transferred to a real wheeled platform. Simulation significantly reduces training time while preserving real-world success rate.

\textbf{Mapless navigation for ASVs (Autonomous Surface Vehicles)}: DRL is applied to navigation in challenging non-stationary scenarios where real-world experiments are difficult. \textbf{Formal verification} methods are additionally used to identify critical configurations and guarantee safety properties of the learned policy.




% ========================================
\part{Wheeled Locomotion: Kinematics}
% ========================================

\section{Introduction to Kinematics for Mobile Robots}

\subsection{Wheeled Locomotion: Motivation}

Wheeled locomotion is:
\begin{itemize}
    \item very efficient on flat/hard surfaces
    \item restricted to man-made infrastructures
    \item the de-facto standard for mobile robots
\end{itemize}

\subsection{Mobile Robots vs.\ Manipulator Arms}

\begin{itemize}
    \item \textbf{Robotic arms}: fixed to the ground; single chain of actuated links (usually).
    \item \textbf{Mobile robots}: position is unbounded; subject to rolling and sliding constraints at the wheel-ground contact point.
\end{itemize}

Key consequences for mobile robot kinematics:
\begin{itemize}
    \item encoder values do not map to unique poses — the workspace is unbounded
    \item no direct way to measure pose: positions must be integrated over time (path-dependent)
    \item integration accumulates errors in position/motion estimates
    \item the constraints imposed by wheels on mobility must be understood
\end{itemize}

\subsection{Non-Holonomic Systems}

\begin{tcolorbox}[colback=blue!5!white,colframe=blue!75!black,title=\textbf{Non-Holonomic System}]
A system is \textbf{non-holonomic} if the kinematic constraints cannot be integrated to a finite position constraint. Knowing the travel distance of each wheel is insufficient to determine the final pose — one must know the full trajectory as a function of time.
\end{tcolorbox}

This is in sharp contrast with robotic arms, where encoder positions map directly to the end-effector pose.

\subsection{Forward and Differential Kinematics}

For mobile robots, the relevant kinematic problems are:

\begin{itemize}
    \item \textbf{Forward kinematics}: given actuator positions (or speeds), determine the corresponding pose (or body velocity).
    \item \textbf{Inverse kinematics}: given a desired pose, determine the corresponding actuator positions — required for motion control.
    \item \textbf{Differential forward kinematics}: given actuator speeds $\dot{\varphi}_i$, determine body velocity $\dot{\xi}$.
    \item \textbf{Differential inverse kinematics}: given a desired body velocity, determine the required actuator speeds.
\end{itemize}

Because non-holonomic constraints prevent direct integration, wheels provide only \textit{differential} inverse kinematics:
\[
\begin{bmatrix} \dot{x} \\ \dot{y} \end{bmatrix} = \begin{bmatrix} \dot{\varphi} r \\ 0 \end{bmatrix}
\]

\section{Wheel Arrangement}

\subsection{Three Design Concepts}

The choice of wheel configuration involves three competing properties:
\begin{itemize}
    \item \textbf{Stability} (static configuration): three wheels are sufficient to guarantee stability, as the center of gravity lies within their support triangle. More than three wheels require a flexible suspension system (hyperstatic configuration).
    \item \textbf{Maneuverability}: how the robot can move in the environment.
    \item \textbf{Controllability}: how actuators must be commanded to achieve motion.
\end{itemize}

No single configuration maximizes all three — the selection depends on the application.

\subsection{Basic Wheel Types}

\begin{itemize}
    \item \textbf{Standard (fixed) wheel} — 2 DOF: rotation around the (motorized) wheel axle and the contact point.
    \item \textbf{Castor wheel} — 3 DOF: rotation around the wheel axle, the contact point, and the castor axle (offset pivot).
    \item \textbf{Swedish (Mecanum) wheel} — 3 DOF: rotation around the wheel axle, the contact point, and around the passive rollers mounted on the rim. Allows motion in any direction.
    \item \textbf{Spherical wheel} — 3 DOF: rotation around the three orthogonal axes.
\end{itemize}

\subsection{Wheel Configurations and Stability}

\begin{itemize}
    \item Stability guaranteed with three wheels: center of gravity (CoG) must lie within the triangle formed by the ground contact points.
    \item Four or more wheels improve stability but create a hyperstatic arrangement requiring flexible suspension.
    \item Bigger wheels overcome higher obstacles, at the cost of higher torque or gear reduction (increased complexity).
    \item Most arrangements are non-holonomic $\Rightarrow$ higher control effort.
    \item Actuation and steering on the same wheel leads to complex design and additional odometry errors.
\end{itemize}

Common layouts include 2-, 3-, 4-, and 6-wheel configurations. A notable example is the \textbf{CMU Uranus}: four independent Swedish wheels at 45°, enabling motion along any trajectory on the plane and simultaneous spinning about the vertical axis.

\section{Mobile Robot Kinematics}

\subsection{Representing Robot Position}

The robot chassis is modeled as a rigid body on a planar environment. Two reference frames are used:
\begin{itemize}
    \item \textbf{Global frame} $\{x_I, y_I\}$: inertial (world) frame.
    \item \textbf{Local frame} $\{x_R, y_R\}$: robot-fixed frame.
\end{itemize}

The robot pose is described by a reference point $P$ with 3 DOF:
\[
\xi_I = \begin{bmatrix} x \\ y \\ \theta \end{bmatrix}
\]
where $(x, y)$ is the position of $P$ in the global frame and $\theta$ is the heading angle.

\subsection{Frame Mapping: Rotation Matrix}

Orthogonal rotation matrices map between local and global frames:
\[
R(\theta) = \begin{bmatrix} \cos\theta & \sin\theta & 0 \\ -\sin\theta & \cos\theta & 0 \\ 0 & 0 & 1 \end{bmatrix}
\]

The velocity in the robot frame relates to the global velocity by:
\[
\dot{\xi}_R = R(\theta)\,\dot{\xi}_I
\]

\subsection{Forward Kinematic Model (Differential Drive)}

For a differential drive robot with wheel radius $r$, axle half-length $l$, wheel spinning speeds $\dot{\varphi}_1$ (right) and $\dot{\varphi}_2$ (left):

\begin{tcolorbox}[colback=blue!5!white,colframe=blue!75!black,title=\textbf{Forward Kinematics: Differential Drive}]
\textbf{Local frame velocity}:
\[
\begin{bmatrix} \dot{x}_R \\ \dot{y}_R \\ \dot{\theta}_R \end{bmatrix}
= \begin{bmatrix} \dfrac{r\dot{\varphi}_1}{2} + \dfrac{r\dot{\varphi}_2}{2} \\[6pt] 0 \\[6pt] \dfrac{r\dot{\varphi}_1}{2l} - \dfrac{r\dot{\varphi}_2}{2l} \end{bmatrix}
\]

\textbf{Global frame velocity} (via $\dot{\xi}_I = R(\theta)^{-1}\dot{\xi}_R$):
\[
\begin{bmatrix} \dot{x}_I \\ \dot{y}_I \\ \dot{\theta}_I \end{bmatrix}
= \begin{bmatrix} \cos\theta & -\sin\theta & 0 \\ \sin\theta & \cos\theta & 0 \\ 0 & 0 & 1 \end{bmatrix}
\begin{bmatrix} \dfrac{r\dot{\varphi}_1}{2} + \dfrac{r\dot{\varphi}_2}{2} \\[6pt] 0 \\[6pt] \dfrac{r\dot{\varphi}_1}{2l} - \dfrac{r\dot{\varphi}_2}{2l} \end{bmatrix}
\]
\end{tcolorbox}

\subsection{Wheel Kinematic Constraints}

Individual wheel constraints are derived under assumptions: wheel in vertical plane, single contact point, no sliding. Two constraint types are imposed for each wheel:

\subsubsection*{Fixed Standard Wheel}

Parameters: $\alpha$ (wheel orientation w.r.t.\ robot frame), $\beta$ (steering angle, fixed here), $l$ (distance from $P$), $r$ (radius).

\begin{itemize}
    \item \textbf{Rolling constraint}: all motion in the wheel plane must be due to rotation,
    \[
    [\sin(\alpha+\beta) \; {-\cos(\alpha+\beta)} \; {(-l)\cos\beta}]\,R(\theta)\dot{\xi}_I - r\dot{\varphi} = 0
    \]
    \item \textbf{Sliding (no-slip) constraint}: no motion perpendicular to wheel plane,
    \[
    [\cos(\alpha+\beta) \; \sin(\alpha+\beta) \; l\sin\beta]\,R(\theta)\dot{\xi}_I = 0
    \]
\end{itemize}

\subsubsection*{Swedish Wheel}

Rolling constraint is modified by the roller angle $\gamma$; there is no constraint on the orthogonal direction (passive rollers absorb lateral motion):
\[
[\sin(\alpha+\beta+\gamma) \; {-\cos(\alpha+\beta+\gamma)} \; {(-l)\cos(\beta+\gamma)}]\,R(\theta)\dot{\xi}_I - r\dot{\varphi}\cos\gamma = 0
\]

\subsubsection*{Complete Robot Model}

For a robot with $M = N_f + N_s$ wheels ($N_f$ fixed, $N_s$ steerable), only fixed and steerable standard wheels impose constraints. Stacking all constraints:

\[
J_1(\beta_s)\,R(\theta)\,\dot{\xi}_I - J_2\,\dot{\varphi} = 0 \qquad \text{(rolling)}
\]
\[
C_1(\beta_s)\,R(\theta)\,\dot{\xi}_I = 0 \qquad \text{(no lateral slip)}
\]

where $J_2 = \mathrm{diag}(r_1, \ldots, r_M)$ and $J_1$, $C_1$ encode the geometric parameters of each wheel.

\section{Mobile Robot Maneuverability}

\subsection{Degrees of Mobility and Steerability}

\begin{tcolorbox}[colback=blue!5!white,colframe=blue!75!black,title=\textbf{Maneuverability}]
Robot maneuverability $\delta_M = \delta_m + \delta_s$ where:
\begin{itemize}
    \item $\delta_m$: \textbf{degree of mobility} — DoF that can be directly manipulated through changes in wheel velocity (from the sliding constraint null space).
    \item $\delta_s$: \textbf{degree of steerability} — number of independently controllable steering parameters.
\end{itemize}
\end{tcolorbox}

\subsection{Five Basic Three-Wheel Configurations}

\begin{center}
\begin{tabular}{lccc}
\hline
Configuration & $\delta_M$ & $\delta_m$ & $\delta_s$ \\
\hline
Omnidirectional & 3 & 3 & 0 \\
Differential drive & 2 & 2 & 0 \\
Omni-steer & 3 & 2 & 1 \\
Tricycle & 2 & 1 & 1 \\
Two-steer & 3 & 1 & 2 \\
\hline
\end{tabular}
\end{center}

\section{Motion Control}

\subsection{Overview}

The objective of a kinematic controller is to follow a trajectory described by position and/or velocity profiles as a function of time. Motion control is challenging because:
\begin{itemize}
    \item mobile robots are typically non-holonomic
    \item they are Multi-Input Multi-Output (MIMO) systems
\end{itemize}
Most controllers (including those covered here) operate at the kinematic level and do not account for system dynamics.

\subsection{Open Loop Control}

Given a trajectory, decompose it into \textbf{motion segments} of clearly defined shape (straight lines or circular arcs). Pre-compute a smooth trajectory for each segment.

Drawbacks:
\begin{itemize}
    \item pre-computing trajectories that incorporate robot constraints is non-trivial
    \item dynamic changes in the environment cannot be handled
    \item resulting overall trajectory may not be smooth at segment transitions
\end{itemize}

\subsection{Feedback Control}

Find a control matrix $K$ (with $k_{ij} = k(t, e)$) such that:
\[
\begin{bmatrix} v(t) \\ \omega(t) \end{bmatrix} = K \cdot e = K \cdot \begin{bmatrix} x \\ y \\ \theta \end{bmatrix}
\]
drives the error to zero: $\lim_{t \to \infty} e(t) = 0$.

\subsubsection*{Kinematic Position Control for Differential Drive}

The kinematics of a differential drive in the inertial frame are:
\[
\begin{bmatrix} \dot{x} \\ \dot{y} \\ \dot{\theta} \end{bmatrix}
= \begin{bmatrix} \cos\theta & 0 \\ \sin\theta & 0 \\ 0 & 1 \end{bmatrix}
\begin{bmatrix} v \\ \omega \end{bmatrix}
\]

Assume w.l.o.g.\ the goal is at the origin of the inertial frame. Let $\alpha$ denote the angle between the robot's $x_R$ axis and the vector from the wheel axle center to the goal.

\subsubsection*{Polar Coordinate Transformation}

Transform to polar coordinates centered at the goal:
\[
\rho = \sqrt{\Delta x^2 + \Delta y^2}, \qquad \alpha = -\theta + \mathrm{atan2}(\Delta y, \Delta x), \qquad \beta = -\theta - \alpha
\]

The system dynamics in polar coordinates (for $\alpha \in I_1 = (-\frac{\pi}{2}, \frac{\pi}{2})$):
\[
\begin{bmatrix} \dot{\rho} \\ \dot{\alpha} \\ \dot{\beta} \end{bmatrix}
= \begin{bmatrix} -\cos\alpha & 0 \\ \frac{\sin\alpha}{\rho} & -1 \\ -\frac{\sin\alpha}{\rho} & 0 \end{bmatrix}
\begin{bmatrix} v \\ \omega \end{bmatrix}
\]
For $\alpha \in I_2 = (-\pi, -\frac{\pi}{2}] \cup [\frac{\pi}{2}, \pi]$, use $v = -v$ (backward motion).

\begin{remark}
The transformation is undefined at $x = y = 0$ (the goal itself). For $\alpha \in I_1$ the forward direction points toward the goal; for $\alpha \in I_2$ the backward direction does. One can always initialize so that $\alpha \in I_1$, but the control law must ensure it stays there.
\end{remark}

\subsubsection*{Control Law}

\begin{tcolorbox}[colback=blue!5!white,colframe=blue!75!black,title=\textbf{Kinematic Feedback Control Law}]
\[
v = k_\rho \rho, \qquad \omega = k_\alpha \alpha + k_\beta \beta
\]
This drives the robot to the equilibrium $(\rho, \alpha, \beta) = (0, 0, 0)$. Stability conditions for $\alpha \in I_1$:
\begin{itemize}
    \item \textit{Local exponential stability}: $k_\rho > 0$,\ $k_\beta < 0$,\ $k_\alpha - k_\rho > 0$
    \item \textit{Strong stability} (no direction reversal): additionally $k_\alpha + \tfrac{5}{3}k_\beta - \tfrac{2}{\pi}k_\rho > 0$
\end{itemize}
\end{tcolorbox}

This drives the robot to $(\rho, \alpha, \beta) = (0, 0, 0)$. The closed-loop dynamics for $\alpha \in I_1$:
\[
\begin{bmatrix} \dot{\rho} \\ \dot{\alpha} \\ \dot{\beta} \end{bmatrix}
= \begin{bmatrix} -k_\rho \rho \cos\alpha \\ k_\rho \sin\alpha - k_\alpha\alpha - k_\beta\beta \\ -k_\rho \sin\alpha \end{bmatrix}
\]

\textbf{Stability conditions}:
\begin{itemize}
    \item \textit{Local exponential stability}: $k_\rho > 0$, $k_\beta < 0$, $k_\alpha - k_\rho > 0$
    \item \textit{Strong stability} (no direction reversal during approach): additionally $k_\alpha + \tfrac{5}{3}k_\beta - \tfrac{2}{\pi}k_\rho > 0$, which guarantees $\alpha \in I_1$ for all $t$ if $\alpha(0) \in I_1$.
\end{itemize}

The control signal $v$ always has constant sign — the robot either drives forward or backward without ever inverting direction (``parking without reversing'').



% ========================================
\part{Perception: Introduction and Sensors}
% ========================================

\section{Introduction to Perception}

\subsection{Perception in the Control Loop}

Perception is the subsystem responsible for extracting meaningful information from raw sensor data and feeding it to higher-level reasoning. In the general control scheme, the \textbf{perception pipeline} comprises three stages:

\begin{enumerate}
    \item \textbf{Sensing}: raw data acquisition from the physical environment.
    \item \textbf{Information extraction and interpretation}: feature detection, object recognition, scene understanding.
    \item \textbf{Localization and map building}: estimating robot pose and constructing an environment model.
\end{enumerate}

The output of this pipeline (position estimate + local map) feeds the cognition block, closing the sense--act loop.

\subsection{Perception as Information Compression}

Perception can be viewed as a hierarchy of increasing abstraction, each level compressing information relative to the one below:

\begin{center}
\begin{tabular}{ll}
\hline
\textbf{Level} & \textbf{Examples} \\
\hline
Raw data & Vision, laser, sound \\
Features & Corners, lines, colors \\
Objects & Doors, pedestrians, bicycles \\
Places/Situations & Kitchen, car overtaking \\
\hline
\end{tabular}
\end{center}

Moving up the hierarchy requires increasingly powerful models and semantic reasoning: from geometric \textit{models} for features, to \textit{models/semantics} for objects, to \textit{reasoning about relationships} for situations.

\subsection{Challenges}

Perception is extremely challenging in practice due to:
\begin{itemize}
    \item \textbf{Noise}: all sensors are subject to measurement noise and quantization error.
    \item \textbf{Partial observation}: sensors have limited field of view, range, and resolution.
    \item \textbf{Context dependency}: interpreting raw data often requires situational context that is difficult to encode algorithmically.
\end{itemize}

Humans excel at perception precisely because of common-sense knowledge, something that remains difficult to replicate in artificial systems.

\section{Sensor Classification}

\subsection{What They Measure}

\begin{itemize}
    \item \textbf{Proprioceptive sensors}: measure quantities \textit{internal} to the robot — e.g., motor speed, battery status, joint angles, wheel odometry.
    \item \textbf{Exteroceptive sensors}: measure quantities from the \textit{environment} — e.g., distance to obstacles, light intensity, object positions.
\end{itemize}

\subsection{How They Measure}

\begin{itemize}
    \item \textbf{Passive sensors}: measure energy naturally present in the environment (e.g., cameras detecting ambient light). Their output is highly dependent on environmental conditions.
    \item \textbf{Active sensors}: emit energy and measure the reaction from the environment (e.g., lidar emitting laser pulses, ultrasonic sensors emitting sound). Generally higher performance, but may interfere with each other or the environment.
\end{itemize}

\subsection{Sensor Classification Tables}

The following tables categorize common sensors by function, whether they are proprioceptive (PC) or exteroceptive (EC), and whether they are active (A) or passive (P).

\begin{center}
\begin{tabular}{llcc}
\hline
\textbf{Category} & \textbf{Sensor} & \textbf{PC/EC} & \textbf{A/P} \\
\hline
\multirow{3}{*}{\parbox{3cm}{Tactile sensors\\(contact, proximity)}} & Contact switches, bumpers & EC & P \\
 & Optical barriers & EC & A \\
 & Noncontact proximity sensors & EC & A \\
\hline
\multirow{7}{*}{\parbox{3cm}{\raggedright Wheel/motor sensors\\(speed, position)}} & Brush encoders & PC & P \\
 & Potentiometers & PC & P \\
 & Synchros, resolvers & PC & A \\
 & Optical encoders & PC & A \\
 & Magnetic encoders & PC & A \\
 & Inductive encoders & PC & A \\
 & Capacitive encoders & PC & A \\
\hline
\multirow{3}{*}{\parbox{3cm}{Heading sensors\\(orientation)}} & Compass & EC & P \\
 & Gyroscopes & PC & P \\
 & Inclinometers & EC & A/P \\
\hline
\multirow{4}{*}{\parbox{3cm}{\raggedright Ground-based beacons\\(localization)}} & GPS & EC & A \\
 & Active optical or RF beacons & EC & A \\
 & Active ultrasonic beacons & EC & A \\
 & Reflective beacons & EC & A \\
\hline
\multirow{5}{*}{\parbox{3cm}{Active ranging\\(range, geometry)}} & Reflectivity sensors & EC & A \\
 & Ultrasonic sensor & EC & A \\
 & Laser rangefinder & EC & A \\
 & Optical triangulation (1D) & EC & A \\
 & Structured light (2D) & EC & A \\
\hline
\multirow{2}{*}{\parbox{3cm}{\raggedright Motion/speed sensors}} & Doppler radar & EC & A \\
 & Doppler sound & EC & A \\
\hline
\multirow{3}{*}{\parbox{3cm}{Vision-based sensors}} & CCD/CMOS camera(s) & EC & P \\
 & Visual ranging packages & EC & P \\
 & Object tracking packages & EC & P \\
\hline
\end{tabular}
\end{center}

\section{Individual Sensor Types}

\subsection{Encoders}

An encoder is an electro-mechanical device that converts position (linear or angular) to an analogue or digital signal. In robotics, \textbf{optical incremental encoders} are the standard:

\begin{itemize}
    \item A disk with alternating opaque and transparent segments is attached to the motor shaft.
    \item Two photo-detectors (channels A and B) are offset by a quarter period (\textbf{quadrature encoding}).
    \item The four-state sequence of (A, B) $\in \{(\text{high},\text{low}),\, (\text{high},\text{high}),\, (\text{low},\text{high}),\, (\text{low},\text{low})\}$ encodes both displacement and direction.
    \item Resolution is typically specified in pulses per revolution (PPR); quadrature decoding multiplies effective resolution by 4.
\end{itemize}

Encoders are proprioceptive and passive — they impose no load on the environment and require no emitted energy.

\subsection{Heading Sensors}

Heading sensors determine the robot's \textbf{orientation} (yaw) and inclination (roll, pitch) with respect to a reference frame.

\begin{itemize}
    \item \textbf{Compass} (exteroceptive, passive): senses the absolute direction of Earth's magnetic field. Gives yaw directly but is susceptible to ferromagnetic interference.
    \item \textbf{Gyroscope} (proprioceptive, passive): senses \textit{relative} orientation change. When fused with velocity information, enables \textbf{dead reckoning} (deduced reckoning) — integrating angular rate over time to estimate heading.
    \item \textbf{Inclinometer}: measures static tilt relative to gravity via a damped pendulum or fluid.
\end{itemize}

\subsubsection*{Gyroscopes}

Two main technologies:

\begin{itemize}
    \item \textbf{Mechanical}: exploits the angular momentum of a rapidly spinning wheel. Force applied on one axis generates a reaction (precession) on the orthogonal axis. High accuracy but high cost and significant drift over time.
    \item \textbf{Optical (fiber-optic / ring laser)}: two laser beams travel in opposite directions around a closed optical path. The beam counter-propagating to the rotation direction travels a slightly shorter path (Sagnac effect), producing a measurable phase shift proportional to angular velocity. Lower drift than mechanical; no moving parts.
\end{itemize}

\subsection{Accelerometers}

An accelerometer measures all external forces acting on the robot, including gravity. The working principle is a \textbf{spring-mass-damper} system: a proof mass $m$ attached to a spring $k$ and damper $c$ deflects under acceleration; the deflection is measured capacitively or piezoelectrically.

Modern accelerometers are fabricated as \textbf{MEMS} (Micro Electro-Mechanical Systems), enabling miniaturization and mass production (applications: airbags, game controllers, smartphones).

\begin{remark}
Gravity must be subtracted before integrating accelerometer data to obtain velocity and position. Without careful calibration, double integration accumulates large errors.
\end{remark}

\subsection{Inertial Measurement Units (IMU)}

An IMU integrates gyroscopes and accelerometers (and optionally magnetometers) to jointly estimate:
\begin{itemize}
    \item relative position $(x, y, z)$
    \item orientation: roll, pitch, yaw
    \item velocity and acceleration
\end{itemize}
all with respect to an inertial reference frame. The processing pipeline is:

\[
\underbrace{\text{Rate gyroscope}}_{\dot{\theta}} \xrightarrow{\int} \text{Orientation}
\qquad
\underbrace{\text{Accelerometer}}_{\mathbf{a}_\text{body}} \xrightarrow{R^{-1}} \mathbf{a}_\text{nav} \xrightarrow{-g} \mathbf{a}_\text{linear} \xrightarrow{\int} \mathbf{v} \xrightarrow{\int} \mathbf{p}
\]

Key issues:
\begin{itemize}
    \item Gravity must be compensated in the navigation frame.
    \item Initial velocity must be known.
    \item \textbf{Drift} is the dominant source of error: integrating small biases accumulates unbounded error over time. IMUs must be periodically re-calibrated or fused with absolute sensors (GPS, cameras).
\end{itemize}

\subsection{Beacons and GPS}

Beacons are artificial landmarks that provide point references to compute the position of a mobile platform.

\subsubsection*{Outdoor: GPS}

The Global Positioning System (GPS) consists of a constellation of satellites broadcasting timing signals. A receiver triangulates its position using signals from at least 4 satellites (the fourth satellite is needed to solve for the unknown receiver clock offset).

\begin{itemize}
    \item Standard civilian GPS accuracy: 5--20 meters.
    \item Differential GPS (DGPS): a fixed reference station with known position broadcasts correction signals $\to$ centimeter-level accuracy achievable.
    \item Critical limitation: \textbf{occlusions} (indoors, urban canyons, tunnels) break satellite visibility.
\end{itemize}

\subsubsection*{Indoor: RFID, WiFi, Bluetooth, UWB}

Indoor localization systems are based on the same triangulation/trilateration principles as GPS but use different signal technologies:
\begin{itemize}
    \item Most technologies are sensitive to multipath propagation and interference.
    \item \textbf{UWB} (Ultra-Wideband): short-pulse ToF ranging with centimeter-level precision; more robust to multipath than narrowband systems.
\end{itemize}

\subsection{Range Sensors}

Range sensors are key sensors for localization, mapping, and obstacle detection. They are popular in robotics due to their low cost and ease of interpretation. Two families:

\subsubsection*{Time-of-Flight (ToF)}

Principle: emit a signal, measure the round-trip travel time $t$, compute distance as $d = c \cdot t$.

\begin{itemize}
    \item \textbf{Sonars}: use sound ($c \approx 0.3\,\text{m/ms}$). At 3 m, the echo returns in $\approx 10\,\text{ms}$ — easy to measure. Wide beam opening angle is a limitation (poor angular resolution).
    \item \textbf{Lidars}: use light ($c \approx 0.3\,\text{m/ns}$). At 3 m, the round-trip takes $\approx 10\,\text{ns}$ — measuring nanosecond intervals requires dedicated electronics. High angular resolution and range accuracy. Used for 2D scan profiles or 3D point clouds.
    \item \textbf{Time-of-Flight cameras}: per-pixel ToF measurement using modulated infrared illumination, producing dense depth images.
\end{itemize}

\begin{tcolorbox}[colback=blue!5!white,colframe=blue!75!black,title=\textbf{ToF Accuracy Factors}]
Quality of ToF sensors depends on:
\begin{itemize}
    \item \textbf{Inaccuracies in time measurement} (dominant for lidars)
    \item \textbf{Opening angle of the beam} (dominant for sonars — poor spatial resolution)
    \item \textbf{Target interactions}: reflectivity, opacity, surface angle (specular vs.\ diffuse reflection)
\end{itemize}
Speed of sound $\approx 0.3\,\text{m/ms}$; speed of light $\approx 0.3\,\text{m/ns}$. EM signals travel $10^6\times$ faster than sound.
\end{tcolorbox}

\subsubsection*{Geometric Active Ranging: Structured Light}

Structured light sensors project a known light pattern (e.g., a line, grid, or coded texture using a laser or projector with a prism) onto the scene. A camera displaced from the emitter observes the projected pattern and detects deformations caused by scene geometry.

\begin{itemize}
    \item The \textbf{correspondence problem} is simplified because the projected pattern is known a priori.
    \item Depth is recovered via \textbf{triangulation}: the displacement of observed pattern features relative to their expected positions encodes range.
    \item Common patterns: single stripe (1D), multiple stripes, coded patterns (binary or phase-shifted).
    \item Modern RGB-D cameras (e.g., Intel RealSense, original Microsoft Kinect) use structured infrared light.
\end{itemize}

\end{document}
